\documentclass[12pt,a4paper]{article}
\usepackage[utf8]{inputenc}
\usepackage[T1]{fontenc}
\usepackage[spanish]{babel}
\usepackage{geometry}
\usepackage{xcolor}
\usepackage{listings}
\usepackage{booktabs}
\usepackage{longtable}
\usepackage{hyperref}
\usepackage{fancyhdr}
\usepackage{titlesec}
\usepackage{enumitem}
\usepackage{float}
\usepackage{array}
\usepackage{tcolorbox}
\usepackage{graphicx}

\geometry{margin=2.5cm}

\definecolor{primary}{RGB}{31,73,125}
\definecolor{secondary}{RGB}{68,114,196}
\definecolor{codebg}{RGB}{245,245,245}
\definecolor{codeframe}{RGB}{200,200,200}
\definecolor{warningbg}{RGB}{255,243,205}
\definecolor{successbg}{RGB}{212,237,218}
\definecolor{tipbg}{RGB}{214,234,248}

\lstset{
    backgroundcolor=\color{codebg},
    frame=single,
    rulecolor=\color{codeframe},
    basicstyle=\ttfamily\small,
    keywordstyle=\color{primary}\bfseries,
    commentstyle=\color{gray}\itshape,
    stringstyle=\color{secondary},
    breaklines=true,
    showstringspaces=false,
    language=bash,
}

\pagestyle{fancy}
\fancyhf{}
\fancyhead[L]{\textcolor{primary}{\textbf{Manual de Usuario}}}
\fancyhead[R]{\textcolor{primary}{GeoMIP — k-Particiones}}
\fancyfoot[C]{\thepage}
\fancyfoot[L]{\small Universidad de Caldas}
\fancyfoot[R]{\small 2026}

\titleformat{\section}{\Large\bfseries\color{primary}}{}{0em}{\thesection.\quad}
\titleformat{\subsection}{\large\bfseries\color{secondary}}{}{0em}{\thesubsection.\quad}
\titleformat{\subsubsection}{\normalsize\bfseries}{}{0em}{\thesubsubsection.\quad}

\hypersetup{colorlinks=true, linkcolor=primary, urlcolor=secondary}

\tcbuselibrary{skins}
\newtcolorbox{notebox}[1][]{colback=warningbg,colframe=orange!70!black,
    title={\textbf{Importante}},fonttitle=\small,#1}
\newtcolorbox{tipbox}[1][]{colback=tipbg,colframe=blue!60!black,
    title={\textbf{Consejo}},fonttitle=\small,#1}
\newtcolorbox{successbox}[1][]{colback=successbg,colframe=green!50!black,
    title={\textbf{Resultado esperado}},fonttitle=\small,#1}

\begin{document}

% ─────────────────────────── PORTADA ───────────────────────────
\begin{titlepage}
    \centering
    \vspace*{1cm}
    {\Huge\bfseries\color{primary} Manual de Usuario\\[0.5em]
    \Large Sistema GeoMIP\\[0.3em]
    \large Algoritmo k-MIP Geométrico para k-Particiones}
    \vspace{1.5cm}

    \rule{\linewidth}{2pt}
    \vspace{0.5cm}

    {\large\textbf{Proyecto de Análisis y Diseño de Algoritmos}}\\[0.3em]
    {\large Ingeniería en Sistemas y Computación}\\[0.3em]
    {\large Facultad de Ingeniería — Universidad de Caldas}
    \vspace{1cm}

    \rule{\linewidth}{0.5pt}
    \vspace{0.8cm}

    {\large\textbf{Autores:}}\\[0.5em]
    {\large Jorge Ivan Garcia Torres}\\[0.3em]
    {\large Camilo Botero Merchor}
    \vspace{1cm}

    \rule{\linewidth}{0.5pt}
    \vspace{0.5cm}
    {\large Semestre 2026-1}

    \vfill
    {\small\color{gray} proyecto-analisis-kpartitions / GeoMIP}
\end{titlepage}

\tableofcontents
\newpage

% ─────────────────────────── 1. INTRODUCCIÓN ───────────────────────────
\section{Introducción}

Este manual describe paso a paso cómo usar el sistema GeoMIP para llenar el Excel del profesor con los resultados del algoritmo k-MIP Geométrico y generar el reporte de benchmark comparativo.

\textbf{¿Qué hace el sistema?}
\begin{enumerate}
    \item Lee 249 casos de prueba del Excel del profesor (sistemas de 10 a 25 nodos).
    \item Para cada caso, ejecuta el algoritmo geométrico que encuentra la partición de $k$ subconjuntos con menor pérdida de información integrada ($\Phi$).
    \item Escribe la partición encontrada, la pérdida $\Phi$ y el tiempo de ejecución en las celdas correspondientes del Excel.
    \item Genera un reporte comparando nuestros resultados con los de QNodes (referencia del profesor).
\end{enumerate}

\textbf{Scripts disponibles:}
\begin{longtable}{lp{9cm}}
    \toprule
    Script & Descripción \\
    \midrule
    \endhead
    \texttt{generar\_tpms.py}     & Genera los archivos de datos (TPM) necesarios. Solo se corre una vez. \\
    \texttt{fill\_excel\_runner.py} & Runner principal. Lee el Excel, corre el algoritmo y escribe resultados. \\
    \texttt{benchmark\_mejorado.py} & Calcula métricas comparativas una vez llenado el Excel. \\
    \bottomrule
\end{longtable}

% ─────────────────────────── 2. PUESTA EN MARCHA ───────────────────────────
\section{Puesta en Marcha}

\subsection{Paso 1 — Activar el entorno virtual}

Antes de cualquier comando, activar el entorno virtual:

\begin{lstlisting}
# Windows (Git Bash)
cd ~/OneDrive/Desktop/Uni/9_Semestre/Analisis/proyecto-analisis-kpartitions
source .venv/Scripts/activate
\end{lstlisting}

Cuando el entorno está activo, el prompt cambia a \texttt{(.venv)}.

Luego navegar al directorio de trabajo:
\begin{lstlisting}
cd GeoMIP/src/Method2_Dynamic_Programming_Reformulation
\end{lstlisting}

\subsection{Paso 2 — Generar los archivos de datos (solo la primera vez)}

\begin{notebox}
Este paso solo es necesario si los archivos \texttt{N20A.csv}, \texttt{N22A.csv} o \texttt{N25A.csv} no existen en la carpeta \texttt{GeoMIP/data/samples/}.
\end{notebox}

Verificar qué archivos existen:
\begin{lstlisting}
ls ../../../GeoMIP/data/samples/
\end{lstlisting}

Generar los faltantes (ejemplo: solo N20A y N22A si N25A es muy grande):
\begin{lstlisting}
python generar_tpms.py --solo 20 22
\end{lstlisting}

Para generar todos (requiere $\sim$4 GB de RAM y varios minutos para N25A):
\begin{lstlisting}
python generar_tpms.py
\end{lstlisting}

\begin{successbox}
Salida esperada para N10A (ya existente):
\begin{verbatim}
  Generando N10A.csv | 1,024 estados x 10 nodos | <1 KB
     Matriz generada en 0.00s -> guardando...
     Guardado en 0.00s
\end{verbatim}
\end{successbox}

% ─────────────────────────── 3. LLENAR EL EXCEL ───────────────────────────
\section{Llenar el Excel del Profesor}

\subsection{Comando básico}

Correr para todos los valores de k (2, 3, 4, 5) sobre los 249 casos:

\begin{lstlisting}
python fill_excel_runner.py \
  --excel "../../../DocsNuevos/DatosPruebas2026_1.xlsx"
\end{lstlisting}

\begin{notebox}
Este comando puede tardar varias horas para casos de N=20--25. Se recomienda empezar con pruebas parciales (ver sección 3.2).
\end{notebox}

\subsection{Prueba rápida — 5 casos por hoja}

Para verificar que todo funciona correctamente antes de la ejecución completa:

\begin{lstlisting}
python fill_excel_runner.py \
  --excel "../../../DocsNuevos/DatosPruebas2026_1.xlsx" \
  --k 2 \
  --limite-por-hoja 5
\end{lstlisting}

Esto ejecuta solo los primeros 5 casos de cada hoja (25 casos en total) usando únicamente $k=2$.

\subsection{Opciones disponibles}

\begin{longtable}{lll}
    \toprule
    Opción & Valor por defecto & Descripción \\
    \midrule
    \endhead
    \texttt{--excel} & (obligatorio) & Ruta al Excel del profesor \\
    \texttt{--k} & 2 3 4 5 & Valores de k a calcular \\
    \texttt{--limite} & ninguno & Límite total de casos \\
    \texttt{--limite-por-hoja} & ninguno & Límite de casos por hoja \\
    \texttt{--timeout} & 300 & Segundos máximos por caso \\
    \texttt{--usar-qnodes} & desactivado & Activa QNodes (experimental) \\
    \bottomrule
\end{longtable}

\subsection{Ejemplos de uso}

\textbf{Solo k=2 y k=3:}
\begin{lstlisting}
python fill_excel_runner.py \
  --excel "../../../DocsNuevos/DatosPruebas2026_1.xlsx" \
  --k 2 3
\end{lstlisting}

\textbf{Con timeout reducido (60 segundos por caso):}
\begin{lstlisting}
python fill_excel_runner.py \
  --excel "../../../DocsNuevos/DatosPruebas2026_1.xlsx" \
  --k 2 \
  --timeout 60
\end{lstlisting}

\textbf{Solo los primeros 10 casos totales:}
\begin{lstlisting}
python fill_excel_runner.py \
  --excel "../../../DocsNuevos/DatosPruebas2026_1.xlsx" \
  --k 2 \
  --limite 10
\end{lstlisting}

\subsection{Interpretar la salida en consola}

Durante la ejecución, cada caso muestra una línea como esta:

\begin{lstlisting}[language={}]
[3/25] 10A-Elementos | fila=8 | condicion=1111111111 |
       alcance=1111111111 | mecanismo=0111111111
CRITICAL: k-MIP geometrico iniciando (k=2).
CRITICAL: Candidatos generados: 13
  Geometric | perdida=0.000000 | tiempo=0.174991s
\end{lstlisting}

\begin{longtable}{lp{9cm}}
    \toprule
    Elemento & Significado \\
    \midrule
    \endhead
    \texttt{[3/25]}         & Caso 3 de 25 totales \\
    \texttt{10A-Elementos}  & Hoja del Excel \\
    \texttt{fila=8}         & Fila del Excel donde se escribirá el resultado \\
    \texttt{condicion=...}  & Cadena binaria de condición (ya convertida de letras) \\
    \texttt{Candidatos: 13} & Número de particiones candidatas evaluadas \\
    \texttt{perdida=0.0}    & $\Phi = 0$: el sistema puede particionarse sin pérdida \\
    \texttt{tiempo=0.17s}   & Tiempo de ejecución del caso \\
    \checkmark \texttt{Geometric} & El caso se resolvió exitosamente \\
    $\times$ \texttt{Geometric} & El caso falló o alcanzó el timeout \\
    \bottomrule
\end{longtable}

\subsection{Resumen final}

Al terminar cada k, se muestra un resumen:

\begin{lstlisting}[language={}]
============================================================
QNodes completados    : 0/25
Geometric completados : 25/25
Fallos registrados    : 0
============================================================
\end{lstlisting}

\begin{tipbox}
Si hay fallos por timeout, se pueden relanzar con \texttt{--timeout} mayor o usar \texttt{--limite-por-hoja} para identificar qué hojas son lentas.
\end{tipbox}

% ─────────────────────────── 4. GENERAR EL BENCHMARK ───────────────────────────
\section{Generar el Reporte de Benchmark}

Una vez llenado el Excel (o al menos parcialmente), se puede generar el reporte de comparación:

\begin{lstlisting}
python benchmark_mejorado.py \
  --excel "../../../DocsNuevos/DatosPruebas2026_1.xlsx"
\end{lstlisting}

Esto crea el archivo \texttt{benchmark\_resultado.xlsx} en el directorio actual con dos hojas:
\begin{itemize}
    \item \textbf{Detalle:} una fila por caso con todas las métricas.
    \item \textbf{Resumen\_por\_hoja:} métricas agregadas por hoja y valor de k.
\end{itemize}

\subsection{Opciones del benchmark}

\begin{lstlisting}
python benchmark_mejorado.py \
  --excel "../../../DocsNuevos/DatosPruebas2026_1.xlsx" \
  --k 2 3 \
  --salida "mi_reporte.xlsx"
\end{lstlisting}

\begin{longtable}{lll}
    \toprule
    Opción & Valor por defecto & Descripción \\
    \midrule
    \endhead
    \texttt{--excel}  & (obligatorio)              & Excel ya llenado por el runner \\
    \texttt{--k}      & 2 3 4 5                    & k(s) a evaluar \\
    \texttt{--salida} & benchmark\_resultado.xlsx  & Archivo de salida con métricas \\
    \texttt{--limite} & ninguno                    & Limitar casos para prueba \\
    \bottomrule
\end{longtable}

\subsection{Interpretar la salida del benchmark}

\begin{lstlisting}[language={}]
================================================================
  BENCHMARK - RESUMEN (Geometric vs QNodes referencia)
================================================================
 [10A-Elementos] k=2 | casos=49 | acierto=95.9% |
   err_rel=0.0001 | jaccard=0.0023 | speedup=42.3x | Excelente
 [15B-Elementos] k=2 | casos=50 | acierto=88.0% |
   err_rel=0.0043 | jaccard=0.0410 | speedup=18.7x | Bueno
================================================================
\end{lstlisting}

\begin{longtable}{lp{9cm}}
    \toprule
    Indicador & Interpretación \\
    \midrule
    \endhead
    🟢 Excelente & Acierto $>90\%$, error $<1\%$, Jaccard $<0.1$ \\
    🟡 Bueno     & Acierto $>80\%$, error $<5\%$, Jaccard $<0.2$ \\
    🟠 Aceptable & Acierto $>70\%$, error $<10\%$, Jaccard $<0.3$ \\
    🔴 Insuficiente & No alcanza los umbrales anteriores \\
    Speedup & Cuántas veces más rápido es Geometric vs QNodes \\
    \bottomrule
\end{longtable}

% ─────────────────────────── 5. FLUJO COMPLETO ───────────────────────────
\section{Flujo de Trabajo Completo}

A continuación se muestra el flujo recomendado de principio a fin:

\begin{lstlisting}
# 0. Activar entorno y navegar
source .venv/Scripts/activate
cd GeoMIP/src/Method2_Dynamic_Programming_Reformulation

# 1. Generar TPMs faltantes (solo primera vez)
python generar_tpms.py --solo 20 22 25

# 2. Prueba rapida: 5 casos por hoja, solo k=2
python fill_excel_runner.py \
  --excel "../../../DocsNuevos/DatosPruebas2026_1.xlsx" \
  --k 2 --limite-por-hoja 5

# 3. Si la prueba funciona, correr completo
python fill_excel_runner.py \
  --excel "../../../DocsNuevos/DatosPruebas2026_1.xlsx"

# 4. Generar reporte comparativo
python benchmark_mejorado.py \
  --excel "../../../DocsNuevos/DatosPruebas2026_1.xlsx" \
  --salida benchmark_resultado.xlsx
\end{lstlisting}

% ─────────────────────────── 6. PROBLEMAS FRECUENTES ───────────────────────────
\section{Solución de Problemas Frecuentes}

\begin{longtable}{p{5cm} p{8cm}}
    \toprule
    Problema & Solución \\
    \midrule
    \endhead
    \texttt{error: unrecognized arguments: --limite-por-hoja} &
    El archivo \texttt{fill\_excel\_runner.py} no está actualizado. Reemplazarlo con la versión más reciente. \\
    \addlinespace
    \texttt{FileNotFoundError: N20A.csv} &
    Ejecutar \texttt{python generar\_tpms.py --solo 20} para generarlo. \\
    \addlinespace
    \texttt{KeyError: Worksheet 25A-Elementos} &
    Error por espacio en el nombre de la hoja. Actualizar \texttt{fill\_excel\_runner.py}. \\
    \addlinespace
    Ejecución muy lenta en N=15+ &
    Usar \texttt{--timeout 60} para limitar a 1 minuto por caso. Los casos que superen el límite se marcan como fallo pero no bloquean los siguientes. \\
    \addlinespace
    El Excel no se actualiza &
    Verificar que el archivo no esté abierto en Excel. Cerrarlo antes de correr el runner. \\
    \addlinespace
    \texttt{(.venv)} no aparece en el prompt &
    El entorno virtual no está activo. Ejecutar \texttt{source .venv/Scripts/activate} desde la raíz del proyecto. \\
    \addlinespace
    Pérdida = 0 en todos los casos &
    Es el resultado correcto para este sistema. $\Phi = 0$ indica independencia causal completa en esas combinaciones de alcance/mecanismo. \\
    \bottomrule
\end{longtable}

% ─────────────────────────── 7. ESTRUCTURA DE ARCHIVOS ───────────────────────────
\section{Estructura de Archivos del Proyecto}

\begin{lstlisting}[language={}]
proyecto-analisis-kpartitions/
|
+-- GeoMIP/
|   +-- data/
|   |   +-- samples/
|   |       +-- N10A.csv     <- Sistema 10 nodos (incluido)
|   |       +-- N15B.csv     <- Sistema 15 nodos (incluido)
|   |       +-- N20A.csv     <- GENERAR: python generar_tpms.py --solo 20
|   |       +-- N22A.csv     <- GENERAR: python generar_tpms.py --solo 22
|   |       +-- N25A.csv     <- GENERAR: python generar_tpms.py --solo 25
|   |
|   +-- src/
|       +-- Method2_Dynamic_Programming_Reformulation/
|           +-- fill_excel_runner.py    <- Runner principal
|           +-- benchmark_mejorado.py   <- Benchmark
|           +-- generar_tpms.py         <- Generador de TPMs
|           +-- benchmark_resultado.xlsx  <- Salida del benchmark
|           |
|           +-- utils/
|           |   +-- excel_loader.py     <- Lectura del Excel
|           |   +-- excel_writer.py     <- Escritura de resultados
|           |
|           +-- src/controllers/
|               +-- manager.py
|               +-- strategies/
|                   +-- geometric_k.py  <- Algoritmo CORE
|                   +-- geometric.py
|                   +-- q_nodes.py
|
+-- DocsNuevos/
    +-- DatosPruebas2026_1.xlsx  <- Excel del profesor (entrada/salida)
\end{lstlisting}

\end{document}